\documentclass[11pt]{article}

\usepackage[margin=1in]{geometry}
\usepackage{amsmath,amssymb,amsthm}
\usepackage{enumitem}
\usepackage{hyperref}
\usepackage{xcolor}

\newtheorem{definition}{Definition}
\newtheorem{assumption}{Assumption}
\newtheorem{theorem}{Theorem}
\newtheorem{lemma}{Lemma}
\newtheorem{corollary}{Corollary}
\newtheorem{proposition}{Proposition}

\newcommand{\State}{\mathcal{S}}
\newcommand{\Actions}{\mathcal{U}}
\newcommand{\Obligations}{\mathcal{H}}
\newcommand{\Valid}{\mathsf{Valid}}
\newcommand{\Bad}{\mathsf{Bad}}
\newcommand{\Clean}{\mathsf{Clean}}
\newcommand{\commit}{\mathsf{commit}}
\newcommand{\abort}{\mathsf{abort}}
\newcommand{\Refine}{\mathsf{Refine}}
\newcommand{\LLM}{\pi_\theta}
\newcommand{\pot}{V}

\title{Transactional Structured Reasoning: Contamination Safety and Recovery}
\author{}
\date{}

\begin{document}
\maketitle

\begin{abstract}
We formalize a transactional execution semantics for structured reasoning with
LLM-generated local proposals. The central object is not the model's hidden chain
of thought, but the committed reasoning state used to produce the final answer.
We prove two basic claims. First, under sound validation, faithful rollback,
and answer-extraction soundness, invalid intermediate proposals cannot
contaminate committed reasoning state, and any emitted final answer is
certified by a committed root certificate. Second, treating the LLM as a
stochastic proposal kernel, retries and refinement \emph{can} amplify the
probability of eventually finding a valid local update without increasing
contamination risk under exact validation, but the amount of amplification
is conditional: it is the i.i.d.\ rate $1-(1-p_h)^{b_h}$ only under sample
independence, is upper-bounded by it under shared latent failure modes, and
is recovered in calibrated form via a persistence bound. These claims are
intended as the reasoning-state analogue of transactional no-regression,
adapted from mutable system state to mutable proof/claim state.
\end{abstract}

\section{Setup}

Let a problem instance be represented by a rooted dependency structure
\[
G = (\Obligations, E, r),
\]
where $\Obligations$ is a finite set of reasoning obligations, $E$ is a
dependency relation, and $r \in \Obligations$ is the root obligation whose
solution determines the final answer.

The committed reasoning state at time $t$ is
\[
s_t = (F_t, D_t, K_t, \sigma_t, \Pi_t, B_t),
\]
where $F_t$ is the actionable frontier, $D_t$ is the set of expanded obligations,
$K_t$ is the set of solved obligations, $\sigma_t$ is the map of committed
values or claims, $\Pi_t$ is the set of certificates/provenance records, and
$B_t$ is the remaining retry budget.

A local proposal is an update
\[
u_t \in \Actions.
\]
Depending on the task, $u_t$ may be a numeric local solution, an executable local
expression, an assignment claim, an elimination claim, or a structural expansion.

The LLM is modeled as a stochastic proposal kernel:
\[
U_{h,k} \sim \LLM(\cdot \mid s, h, k),
\]
where $h$ is the current obligation and $k$ is the attempt index. The executor
does not directly trust samples from $\LLM$. Every proposed update is passed
through a validator
\[
g(s,u) \in \{0,1\}.
\]

\begin{definition}[Committed-state transition]
Given state $s$ and proposal $u$, let $s \oplus u$ denote the candidate state
obtained by applying $u$ to a copy of $s$. The committed transition is
\[
T(s,u) =
\begin{cases}
s \oplus u & \text{if } g(s,u)=1 \text{ and } \pot(s \oplus u) < \pot(s),\\
s & \text{otherwise.}
\end{cases}
\]
The rejected case is a rollback: the committed state remains exactly $s$.
\end{definition}

\begin{definition}[Rejection memory]
The executor may also maintain an auxiliary rejection memory $R_t$ containing
records of rejected proposals, validator reasons, and retry counts. This is
operational feedback, not committed reasoning state. It may condition later
LLM proposals, but it is excluded from $\sigma_t$, $\Pi_t$, and final answer
extraction.

On a rejected proposal $u_t$ for obligation $h_t$, the operational context may
change from $(s_t,R_t)$ to
\[
(s_t, R_t \cup \{(h_t,u_t,\mathrm{reason}_t)\}),
\]
or to a bounded/windowed version of that update. The committed projection
remains exactly $s_t$. On an accepted update, $R_t$ may be reset or scoped to
the new committed state, because the entailment context has changed.
\end{definition}

\begin{definition}[Bad and clean states]
Let $\Bad(s)=1$ if committed state $s$ contains at least one invalid committed
claim, value, or certificate. Otherwise $\Bad(s)=0$. Define
\[
\Clean(s) \iff \Bad(s)=0.
\]
The predicate $\Clean$ concerns committed reasoning state only. The LLM may
generate invalid hidden text, invalid drafts, or invalid candidate proposals;
these are not contamination unless they enter $\sigma_t$ or $\Pi_t$. Rejection
memory $R_t$ is also not contamination, provided final answer extraction ignores
$R_t$ and $R_t$ is not treated as a committed fact store.
\end{definition}

\begin{definition}[Progress potential]
$\pot : \State \to \mathbb{N}$ is a nonnegative integer potential measuring
unresolved work. In the current implementation this may be an action-unit
potential such as
\[
\pot(s)=
|\{h \in \Obligations : h \notin K_t\}|
+
|\{h \in \Obligations : \mathrm{Ch}(h)\neq\emptyset, h\notin D_t\}|.
\]
For logic-grid tasks, $\pot$ may instead count unresolved person obligations and
remaining candidate claims.
\end{definition}

\section{Claim I: Contamination Safety}

The first claim is the reasoning-state analogue of transactional safety: invalid
intermediate thoughts cannot poison the state used for the final answer.

\begin{assumption}[Initial clean state]
\label{ass:initial-clean}
\[
\Clean(s_0).
\]
\end{assumption}

\begin{assumption}[Exclusive commit authority]
\label{ass:exclusive-commit}
Only the transactional executor can mutate committed state fields
$(D_t,K_t,\sigma_t,\Pi_t)$. Model samples, prompts, scratchpads, and rejected
candidate states are not directly committed.
\end{assumption}

\begin{assumption}[Faithful rollback]
\label{ass:faithful-rollback}
If a proposal $u$ is rejected, the next committed state equals the previous
committed state:
\[
T(s,u)=s.
\]
\end{assumption}

\begin{assumption}[Validator soundness]
\label{ass:validator-soundness}
For every clean state $s$ and proposal $u$,
\[
\Clean(s) \wedge g(s,u)=1 \implies \Clean(s \oplus u).
\]
For exact validators, this is enforced by checking the proposed local update
against an oracle evaluator, solver, proof checker, or entailment checker.
\end{assumption}

\begin{assumption}[Final answer reads only committed state]
\label{ass:answer-from-commit}
The final answer is computed by a function
\[
A : \State \to \mathcal{Y}
\]
that depends only on committed fields such as $\sigma_t$ and $\Pi_t$, not on
rejected proposals or hidden LLM drafts.
\end{assumption}

\begin{assumption}[Answer-extraction soundness]
\label{ass:answer-extraction}
The executor emits a final answer only when the root obligation $r$ has been
solved, i.e.\ $r\in K_t$, and the answer is computed deterministically from
a committed certificate $\Pi_t(r)$ that the same validator $g$ has accepted
as a witness for $\sigma_t(r)$ being the correct value of $r$. Equivalently,
the validator is invoked on the root certificate before any answer leaves
the executor, and accepted certificates correctly witness their claims under
the validator's soundness guarantee.
\end{assumption}

Assumption~\ref{ass:answer-from-commit} is purely structural: it says rejected
proposals are not read at answer time. Assumption~\ref{ass:answer-extraction}
is what is needed for \emph{correctness} of an emitted answer; without it,
contamination safety guarantees only that the answer was extracted from a
clean committed state, not that the emitted value is right.

\begin{theorem}[Transactional Reasoning Contamination Safety]
\label{thm:contamination-safety}
Assume the initial clean state, exclusive commit authority, faithful
rollback, and validator soundness assumptions
(Assumptions~\ref{ass:initial-clean}--\ref{ass:validator-soundness}). Then
every committed state reachable by the transactional executor is clean:
\[
\forall t \ge 0,\quad \Clean(s_t).
\]
Consequently, rejected or invalid intermediate proposals cannot contaminate
the reasoning state used to compute the final answer.

(Note that Assumption~\ref{ass:answer-from-commit} is needed only to
guarantee the final answer reads exclusively from this clean committed
state, and Assumption~\ref{ass:answer-extraction} is needed only for the
correctness corollary below; the cleanness of $s_t$ itself does not depend
on either.)
\end{theorem}

\begin{proof}
The proof is by induction on $t$.

\textbf{Base case.} By the initial clean-state assumption, $\Clean(s_0)$.

\textbf{Inductive step.} Assume $\Clean(s_t)$. Consider a proposal $u_t$.
There are two cases.

If $u_t$ is rejected, faithful rollback gives
\[
s_{t+1}=T(s_t,u_t)=s_t.
\]
Therefore $\Clean(s_{t+1})$ by the induction hypothesis.

If $u_t$ is accepted, then $g(s_t,u_t)=1$ and
\[
s_{t+1}=s_t\oplus u_t.
\]
By validator soundness and $\Clean(s_t)$, it follows that
\[
\Clean(s_t\oplus u_t),
\]
and hence $\Clean(s_{t+1})$.

Thus $\Clean(s_t)$ holds for all reachable committed states. Since the final
answer function reads only committed state, rejected or invalid intermediate
proposals cannot affect the final answer except by failing to contribute a valid
committed update.
\end{proof}

\begin{corollary}[No poisoning under exact validation]
If the validator is exact, $\Clean(s_0)$, and the final answer depends only on
committed state, then the final answer cannot be poisoned by invalid rejected
thoughts.
\end{corollary}

\begin{corollary}[Conditional correctness of emitted answers]
\label{cor:conditional-correctness}
Under Assumptions~\ref{ass:initial-clean}--\ref{ass:answer-extraction} with exact validation, if the
executor emits a final answer, that answer is correct.
\end{corollary}

\begin{proof}
By Theorem~\ref{thm:contamination-safety} every committed state $s_t$ is clean. By
Assumption~\ref{ass:answer-extraction} an answer is emitted only when
$r\in K_t$ and the validator has accepted the root certificate $\Pi_t(r)$;
exact validation then implies $\Pi_t(r)$ correctly witnesses $\sigma_t(r)$.
The emitted answer is a deterministic function of these committed fields, so
it agrees with the true value of $r$.
\end{proof}

\paragraph{Interpretation.}
Theorem~\ref{thm:contamination-safety} alone does not give correctness; it gives only that whatever lands
in committed state was accepted by the validator on a clean predecessor.
Correctness of an emitted answer requires the additional
Assumption~\ref{ass:answer-extraction}: the same validator must certify the
root before extraction. Even with both assumptions, the executor may still
fail to emit (retry budget exhausted, stall, no valid proposal), so the
guarantee is conditional on emission. The role of Theorem~\ref{thm:contamination-safety} in this picture
is to ensure the predecessor state on which the root certificate is checked
is itself clean, so the validator's local guarantee composes into a global
correctness guarantee.

\section{Progress and Termination}

Contamination safety prevents bad updates from poisoning the answer. A separate
progress condition prevents the system from accepting valid but unhelpful
updates forever.

\begin{lemma}[Monotone progress]
If every accepted proposal satisfies
\[
\pot(s \oplus u) < \pot(s),
\]
and rejected proposals leave $s$ unchanged, then $\pot(s_t)$ strictly decreases
on accepted steps and never increases on rejected steps.
\end{lemma}

\begin{proof}
Accepted steps decrease $\pot$ by the commit rule. Rejected steps produce
$s_{t+1}=s_t$, so $\pot(s_{t+1})=\pot(s_t)$.
\end{proof}

\begin{theorem}[Finite-budget termination]
Assume $\pot$ is a nonnegative integer, accepted steps strictly decrease $\pot$,
and rejected steps consume finite retry budget. Then the executor terminates
with success, stall, budget exhaustion, or a configured max-step cap.
\end{theorem}

\begin{proof}
There can be at most $\pot(s_0)$ accepted steps before the potential reaches
zero, and at most $B_0$ rejected steps before retry budget is exhausted. A
configured max-step cap gives an additional finite bound.
\end{proof}

\section{Claim II: Recovery by Retry}

Safety alone is insufficient. The useful property is that rejected local
failures are non-catastrophic: after rollback, the executor can try again from
the same clean state.

Fix a clean state $s$ and a frontier obligation $h$. Define the one-shot local
success probability
\[
p_h(s) =
\Pr_{U \sim \LLM(\cdot \mid s,h)}
\left[
g(s,U)=1
\ \wedge\
\pot(s\oplus U)<\pot(s)
\right].
\]
This is the probability that one LLM sample proposes an acceptable, progressive
local update.

\begin{theorem}[Retry amplification, i.i.d.\ regime]
\label{thm:retry-amplification}
Suppose the executor attempts obligation $h$ up to $b_h$ times from the same
clean committed state $s$, rejected attempts leave $s$ unchanged, and attempts
are conditionally independent with identical local success probability $p_h(s)$.
Then the probability of at least one valid commit for $h$ is
\[
q_h(s,b_h) =
1-(1-p_h(s))^{b_h}.
\]
\end{theorem}

\begin{proof}
Each attempt fails with probability $1-p_h(s)$. By conditional independence, all
$b_h$ attempts fail with probability
\[
(1-p_h(s))^{b_h}.
\]
Therefore at least one attempt succeeds with probability
\[
1-(1-p_h(s))^{b_h}.
\]
Because rejected attempts roll back to the same clean committed state, failed
samples do not alter the success probability by contaminating future attempts.
\end{proof}

\begin{corollary}[Almost-sure local recovery under nonzero support]
If $p_h(s)>0$, then
\[
\lim_{b_h\to\infty} q_h(s,b_h)=1.
\]
\end{corollary}

\paragraph{Interpretation.}
The LLM need not be reliable globally. It only needs to assign nonzero
probability mass to valid local proposals. Transactions allow the executor to
sample this distribution repeatedly while preventing invalid samples from
corrupting state.

\section{Sample Correlation and the Limits of Retry Amplification}
\label{sec:correlation}

Theorem~\ref{thm:retry-amplification} assumes that the $b_h$ retries on obligation $h$ are conditionally
independent. Empirically, repeated samples from a fixed LLM on a fixed prompt
are correlated: failure modes tend to repeat. The realistic generative model
has a latent ``hardness'' or ``failure-mode'' variable $Z$ that is sampled
once per obligation (encoding, e.g., which sub-pattern of the prompt the model
is currently confusing), with attempts conditionally independent given $Z$.

\begin{proposition}[Correlation makes retry amplification an upper bound]
\label{prop:correlation-upper-bound}
Suppose attempts on obligation $h$ share a latent variable $Z$ such that,
conditional on $Z$, the failure indicators $Y_k=\mathbf{1}[X_k=0]$ are
i.i.d.\ Bernoulli$(\rho(Z))$ for some measurable $\rho:Z\to[0,1]$ with
$\mathbb{E}[\rho(Z)]=1-p_h$. Then
\[
\Pr[Y_1=\cdots=Y_{b_h}=1]
=\mathbb{E}\!\left[\rho(Z)^{b_h}\right]
\;\ge\;
\bigl(\mathbb{E}[\rho(Z)]\bigr)^{b_h}
=(1-p_h)^{b_h},
\]
by Jensen's inequality applied to the convex function $x\mapsto x^{b_h}$.
Equivalently,
\[
\Pr[\text{success in } b_h \text{ tries}]
\;\le\;
1-(1-p_h)^{b_h}.
\]
\end{proposition}

The bound of Theorem~\ref{thm:retry-amplification} is therefore an \emph{upper} bound on the achievable
retry amplification under the latent-mode model, not a guarantee. The gap
$\mathbb{E}[\rho(Z)^{b_h}]-(1-p_h)^{b_h}$ grows with the variance of $\rho(Z)$:
when failure mode is essentially deterministic per-prompt ($\rho(Z)\in\{0,1\}$
almost surely), retries achieve no amplification and the realised success rate
is exactly $p_h$ regardless of $b_h$.

\paragraph{Diversity restores amplification.}
The framework therefore implicitly relies on diversity-inducing machinery to
make samples behave more like independent draws and reduce the variance of
$\rho(Z)$. Practical mechanisms include sampling temperature, prompt
paraphrasing, varied few-shot exemplars, distinct draft variables, and
ensembling across decoders or model instances. Each such intervention can be
viewed as enlarging the support of $Z$ relative to a fixed-prompt baseline.

\paragraph{A more conservative practical bound.}
For empirical use, a pessimistic single-parameter alternative to Theorem~\ref{thm:retry-amplification} is
the conditional persistence bound. Let $F_k=\{Y_1=\cdots=Y_k=1\}$ be the
event that the first $k$ attempts on obligation $h$ all fail, and assume the
\emph{uniform} persistence condition
\[
\Pr[Y_{k+1}=1 \mid F_k] \;\le\; \rho_h
\quad \text{for all } k=1,\dots,b_h-1,
\]
with $\rho_h\in[1-p_h,1]$. This is a stronger assumption than mere one-step
correlation: failures must remain at most $\rho_h$-likely conditional on
\emph{all} previous failures, not only on the immediately previous one. Then,
applying the chain rule,
\[
\Pr[F_{b_h}]
\;=\;
\Pr[Y_1=1]\prod_{k=1}^{b_h-1}\Pr[Y_{k+1}=1\mid F_k]
\;\le\;
(1-p_h)\,\rho_h^{\,b_h-1},
\]
so
\[
\Pr[\text{success in } b_h \text{ tries}]
\;\ge\;
1 - (1-p_h)\,\rho_h^{\,b_h-1}.
\]
The bound interpolates between independence ($\rho_h=1-p_h$, recovering the
Theorem~\ref{thm:retry-amplification} rate) and total correlation ($\rho_h=1$, no amplification beyond
the first attempt). The uniform persistence parameter $\rho_h$ can be
calibrated from logged retry trajectories by computing, for each prefix
length $k$, the empirical fraction of trajectories whose $(k{+}1)$-st
attempt fails given $F_k$, and taking the supremum over $k$.

\section{Global Success Under Oracle Structure}

Let $\mathcal{R} \subseteq \Obligations$ be the set of required obligations for a
successful derivation under a fixed oracle structure, topologically ordered
as $h_1,\ldots,h_M$ along dependency edges. Suppose the executor must obtain
one valid commit for each $h\in\mathcal{R}$. We separate the structural
chain-rule argument from the local-rate model.

\begin{theorem}[Product lower bound, abstract form]
\label{thm:product-abstract}
Assume exact validation and rollback safety. Suppose that for each obligation
$h_i$, conditional on all prior obligations $h_1,\ldots,h_{i-1}$ having been
resolved (i.e.\ on reaching a clean committed state in which $h_i$ becomes
the active frontier element), the executor resolves $h_i$ with probability
at least $q_{h_i}^{\mathrm{LB}}$. Then
\[
\Pr[\mathrm{success}]
\;\ge\;
\prod_{i=1}^{M} q_{h_i}^{\mathrm{LB}}.
\]
\end{theorem}

\begin{proof}
Let $E_i$ denote the event that $h_i$ is resolved. By the chain rule,
\[
\Pr\!\left[\bigcap_{i=1}^{M}E_i\right]
=\prod_{i=1}^{M}\Pr\!\left[E_i\;\bigg|\;\bigcap_{j<i}E_j\right]
\;\ge\;
\prod_{i=1}^{M} q_{h_i}^{\mathrm{LB}},
\]
using the assumed conditional lower bound at each topological position.
Exact validation and clean predecessor states (Theorem~\ref{thm:contamination-safety}) ensure the
preconditions of each $q_{h_i}^{\mathrm{LB}}$ are satisfied along any path
that reaches $h_i$.
\end{proof}

The role of any specific local model — i.i.d.\ retries, persistence bound,
exchangeable mixture — is to supply $q_{h_i}^{\mathrm{LB}}$. Three natural
instantiations:

\begin{corollary}[Instantiation: independent retries]
\label{cor:product-iid}
If, conditional on reaching the state at which $h_i$ is attempted, the
$b_{h_i}$ retries are i.i.d.\ with marginal success probability at least
$p_{h_i}$, then $q_{h_i}^{\mathrm{LB}} = 1-(1-p_{h_i})^{b_{h_i}}$ and
\[
\Pr[\mathrm{success}]
\;\ge\;
\prod_{i=1}^{M}\left(1-(1-p_{h_i})^{b_{h_i}}\right).
\]
\end{corollary}

\begin{corollary}[Instantiation: persistence-calibrated retries]
\label{cor:product-persistence}
If, conditional on reaching the state at which $h_i$ is attempted, retries on
$h_i$ satisfy the uniform persistence condition with parameter
$\rho_{h_i}\in[1-p_{h_i},1]$, then
$q_{h_i}^{\mathrm{LB}} = 1 - (1-p_{h_i})\rho_{h_i}^{\,b_{h_i}-1}$ and
\[
\Pr[\mathrm{success}]
\;\ge\;
\prod_{i=1}^{M}\left(1-(1-p_{h_i})\,\rho_{h_i}^{\,b_{h_i}-1}\right).
\]
\end{corollary}

\begin{corollary}[Instantiation: deterministic per-prompt failure]
\label{cor:product-deterministic}
If on each obligation the failure mode is essentially deterministic given the
prompt ($\rho(Z)\in\{0,1\}$ a.s.\ in the latent-mode model of
Proposition~\ref{prop:correlation-upper-bound}), then
$q_{h_i}^{\mathrm{LB}} = p_{h_i}$ regardless of $b_{h_i}$, and the lower
bound degenerates to the one-shot structured baseline
$\prod_i p_{h_i}$.
\end{corollary}

\begin{corollary}[Uniform i.i.d.\ lower bound]
If, in addition to the i.i.d.\ instantiation, $p_h \ge p_{\min}>0$ and
$b_h\ge b$ for all $h\in\mathcal{R}$, and $|\mathcal{R}|=M$, then
\[
\Pr[\mathrm{success}]
\;\ge\;
\left(1-(1-p_{\min})^b\right)^M.
\]
\end{corollary}

\paragraph{Comparison to non-transactional structured execution.}
A one-shot structured baseline that commits the first local proposal succeeds
with probability approximately $\prod_{h\in\mathcal{R}} p_h$ (and is also
unsafe: rejected proposals can corrupt downstream state). Transactional retry
changes the local factor from $p_h$ to $q_h^{\mathrm{LB}}$, where
$q_h^{\mathrm{LB}}\in[p_h,\, 1-(1-p_h)^{b_h}]$ depending on the diversity of
samples within an obligation. Exact validation simultaneously prevents
failed attempts from contaminating downstream obligations.

\paragraph{Worked example: same $p_h$, four different global rates.}
The choice of local-rate model dominates the global success probability,
even with $p_h$ and $b_h$ held fixed. As an illustration take $M=10$
identical obligations, per-attempt local success rate $p_h=0.30$, retry
budget $b_h=8$, and persistence parameter $\rho_h=0.85$. Then:

\begin{center}
\begin{tabular}{l l l}
\hline
Regime & Per-obligation $q_h^{\mathrm{LB}}$ & Global $\Pr[\text{success}]$ \\
\hline
One-shot baseline ($b_h=1$) & $0.30$ & $\approx 5.9 \times 10^{-6}$ \\
Deterministic per-prompt failure & $0.30$ & $\approx 5.9 \times 10^{-6}$ \\
Persistence-calibrated, $\rho_h=0.85$ & $1-(1-p_h)\rho_h^{\,b_h-1}\approx 0.776$ & $\approx 7.9 \times 10^{-2}$ \\
I.i.d.\ retries & $1-(1-p_h)^{b_h}\approx 0.942$ & $\approx 5.5 \times 10^{-1}$ \\
\hline
\end{tabular}
\end{center}

The i.i.d.\ regime beats the persistence-calibrated regime by roughly a
factor of $7$, and the persistence-calibrated regime beats the deterministic
regime by roughly a factor of $1.3\times 10^4$. Note that the deterministic
regime is indistinguishable from the one-shot baseline at the global level:
when each obligation's failure mode is fixed by the prompt, retries within
an obligation are wasted and only diversity \emph{across} prompts (or
refinement) can help. The practical implication is that the operative
quantity is not $b_h$ but the variance of $\rho(Z)$ that the executor's
diversity machinery induces.

\section{Recovery by Refinement}

Retry is not the only recovery operator. If an obligation is too coarse, the
executor can refine it into smaller obligations.

\begin{definition}[Refinement]
A refinement operator replaces a failed obligation $h$ with child obligations
\[
\Refine(h)=\{c_1,\ldots,c_m\}
\]
plus, optionally, a recomposition obligation that combines the children into a
solution for $h$.
\end{definition}

Let $q_h$ denote the (calibrated or assumed) probability that the executor
resolves $h$ directly under its allocated retry budget — that is, the
$q_h^{\mathrm{LB}}$ supplied by whichever local-rate model
(Corollaries~\ref{cor:product-iid}--\ref{cor:product-deterministic}) applies.
Let $q_{c_i}$ denote the analogous resolution probability for child $c_i$.
Let $q_{\mathrm{exp}}$ be the probability that the refinement step itself is
proposed validly and accepted by the validator, and let $q_{\mathrm{comp}}$
be the probability that the children can be recomposed into a validated
solution for $h$.

\begin{proposition}[When refinement improves recovery]
\label{prop:refinement-helps}
Suppose that, conditional on refinement being accepted, the events
``$c_i$ resolved'' for $i=1,\ldots,m$ and ``recomposition succeeds'' are
mutually independent — a substantive assumption, equivalent to asserting
that no shared latent failure mode couples the children. Then refinement
improves the probability of resolving $h$ whenever
\[
q_{\mathrm{exp}}
\left(\prod_{i=1}^m q_{c_i}\right)
q_{\mathrm{comp}}
\;>\;
q_h.
\]
Without the cross-child independence assumption, the same comparison holds
with $\prod_{i=1}^m q_{c_i}$ replaced by the joint resolution probability
$\Pr[\text{all children resolved}\mid \text{refinement accepted}]$, which
must be supplied by the local-rate model.
\end{proposition}

\begin{proof}
The left-hand side is the probability that refinement is accepted, every
child obligation is resolved, and recomposition succeeds, assuming the
stated factorisation across children. The right-hand side is the calibrated
probability of resolving $h$ without refinement. Refinement is beneficial
exactly when the former exceeds the latter. The fall-back form follows by
replacing the product with the joint probability and applying the same
inequality.
\end{proof}

\paragraph{Interpretation.}
This is a decision rule for adaptive decomposition under whichever model of
local resolution the system has calibrated. When local rates are $p_h$ and
samples are i.i.d., $q_h$ takes the form $1-(1-p_h)^{b_h}$ and the
comparison reduces to the original retry-vs-refinement formula. When samples
are correlated, $q_h$ may be much smaller than this i.i.d.\ value, and the
comparison should use the calibrated $q_h$ rather than substituting the
i.i.d.\ formula. Naively using the i.i.d.\ amplified rate when samples share
a latent failure mode tends to make refinement look worse than it is, since
diverse children can break a shared failure mode that a single coarser
obligation cannot.

\section{Noisy Validators}

Exact validators are the cleanest case. For freeform reasoning, validators may
be noisy. Two distinct false-accept rates that are easy to confuse must be
kept separate:
\[
\alpha_h^{\rightarrow}
=
\Pr[g(s,U)=1 \mid U \text{ is invalid}]
\quad\text{(conditional false-accept rate),}
\]
\[
\alpha_h^{\leftarrow}
=
\Pr[U \text{ is invalid} \mid g(s,U)=1]
\quad\text{(posterior false-accept rate),}
\]
\[
\beta_h
=
\Pr[g(s,U)=0 \mid U \text{ is valid}]
\quad\text{(false-reject rate).}
\]
The two false-accept rates are related by Bayes,
\[
\alpha_h^{\leftarrow}
=
\frac{\alpha_h^{\rightarrow}\,\Pr[U \text{ invalid}]}
     {\alpha_h^{\rightarrow}\,\Pr[U \text{ invalid}]
      +(1-\beta_h)\,\Pr[U \text{ valid}]},
\]
and $\alpha_h^{\leftarrow}$ can be much larger than $\alpha_h^{\rightarrow}$
when the proposal kernel is weak (high $\Pr[U\text{ invalid}]$) or the
validator is conservative on valid proposals (high $\beta_h$). The two yield
genuinely different contamination bounds.

\begin{theorem}[Contamination bound, per-invocation form]
\label{thm:contamination-noisy-invocation}
Let $N$ be the (random) total number of validator invocations during a run,
with deterministic upper bound $\bar{N}$ given by retry budget plus the commit
cap. If every validator invocation has conditional false-accept rate at most
$\alpha_{\max}^{\rightarrow}$, then
\[
\Pr[\exists \text{ bad committed update}]
\;\le\;
\bar{N}\,\alpha_{\max}^{\rightarrow}.
\]
\end{theorem}

\begin{proof}
Index validator invocations by $j=1,\dots,N$. Let $V_j=\mathbf{1}[g_j=1]$ and
$I_j=\mathbf{1}[U_j \text{ invalid}]$. A bad commit at invocation $j$ is the
event $V_j=1\wedge I_j=1$, with
\[
\Pr[V_j=1, I_j=1]
=\Pr[V_j=1\mid I_j=1]\,\Pr[I_j=1]
\le \alpha_{\max}^{\rightarrow}\cdot 1.
\]
Union-bounding over the at most $\bar{N}$ invocations,
\[
\Pr\!\left[\bigcup_{j=1}^{N}\{V_j=1, I_j=1\}\right]
\le \sum_{j=1}^{\bar{N}}\Pr[V_j=1, I_j=1]
\le \bar{N}\,\alpha_{\max}^{\rightarrow}. \qedhere
\]
\end{proof}

\begin{theorem}[Contamination bound, per-commit form]
\label{thm:contamination-noisy-posterior}
If the executor produces at most $K$ accepted commits and every accepted commit
satisfies $\Pr[U\text{ invalid}\mid g(s,U)=1]\le\alpha_{\max}^{\leftarrow}$,
then
\[
\Pr[\exists \text{ bad committed update}]
\;\le\;K\,\alpha_{\max}^{\leftarrow}.
\]
\end{theorem}

\begin{proof}
For commit $i\in\{1,\dots,K\}$ let $B_i$ be the event that commit $i$ is
invalid. By assumption $\Pr[B_i]\le\alpha_{\max}^{\leftarrow}$. Union-bounding,
$\Pr[\bigcup_{i=1}^{K}B_i]\le K\,\alpha_{\max}^{\leftarrow}$.
\end{proof}

\paragraph{Why the distinction matters.}
A naive union bound of the form $K\alpha_{\max}^{\rightarrow}$, mixing a commit
count with a conditional false-accept rate, is incorrect. As a counterexample,
suppose all proposals are invalid ($\Pr[U \text{ invalid}]=1$),
$\alpha^{\rightarrow}=10^{-3}$, the retry budget allows $10^4$ invocations,
and the executor halts after $K=1$ accept. Every accept is then a bad commit,
and
\[
\Pr[\exists \text{ bad commit}]
=1-(1-10^{-3})^{10^4}\approx 1,
\]
which exceeds $K\alpha^{\rightarrow}=10^{-3}$ by three orders of magnitude.
Theorem~\ref{thm:contamination-noisy-invocation} gives the correct value
$10^4\cdot 10^{-3}=10$, which is vacuous but not false; this reflects that the
operating regime is genuinely unsafe.

\paragraph{Which bound to use.}
The conditional rate $\alpha_h^{\rightarrow}$ is what one usually estimates
when calibrating a learned validator (held-out invalid examples passed through
the validator). The per-invocation bound is therefore the directly measurable
one, at the cost of paying for every rejected attempt. The per-commit bound is
tighter when it applies, but $\alpha_h^{\leftarrow}$ depends on the proposal
distribution and is not a property of the validator alone.

\paragraph{Effect of false rejects on recovery.}
False rejects do not contaminate state, but they reduce recovery probability.
If the model proposes a valid local update with probability $p_h$ and the
validator rejects valid proposals with probability $\beta_h$, then the
effective accepted local success probability is
\[
p'_h = p_h(1-\beta_h),
\]
and retry amplification becomes
\[
q'_h = 1-(1-p_h(1-\beta_h))^{b_h}.
\]

\section{Open Problems}
\label{sec:open-problems}

The conditional guarantees above leave the substantive design problems open.
We collect them here.

\paragraph{Constructing sound validators.}
Theorem~\ref{thm:contamination-safety} assumes a sound validator. For most freeform reasoning targets
(natural-language proofs, multi-step plans with side effects, open-ended
program synthesis), constructing such a validator \emph{is} the central
problem. The framework is most directly useful for tasks that admit a cheap,
exact local checker: numeric solvers, executable code with deterministic
oracles, formal proofs in a checkable calculus, and constraint problems with
verifier subroutines. Identifying the boundary of validator-friendly tasks,
and designing decompositions that push freeform problems into validator-
friendly sub-obligations, are the substantive research questions.

\paragraph{Optimal budget allocation across obligations.}
Given a total invocation budget $B$ and required obligations
$\mathcal{R}=\{h_1,\dots,h_M\}$ with estimated local success rates
$\hat p_1,\dots,\hat p_M$, the natural objective derived from the product
lower bound is
\[
\max_{b_1,\dots,b_M\ge 0}
\;
\sum_{i=1}^{M}\log\!\left(1-(1-\hat p_i)^{b_i}\right)
\quad\text{s.t.}\quad
\sum_{i=1}^{M} b_i \le B.
\]
This is a separable concave allocation; its KKT conditions yield a
water-filling rule
\[
\frac{(1-\hat p_i)^{b_i}\log\frac{1}{1-\hat p_i}}
     {1-(1-\hat p_i)^{b_i}}
\;=\;
\lambda
\quad\forall i\in\mathcal{R},
\]
for a common dual $\lambda$ tuned to the budget. With unknown $\hat p_i$, the
problem becomes a structured multi-armed bandit: each obligation is a
Bernoulli arm with unknown success probability, retries are pulls, but pulls
also drain a shared budget and feed into a topological dependency. Whether
standard strategies (UCB, Thompson sampling, or an explore-then-allocate
phase based on the water-filling rule above) admit useful regret guarantees
against the offline product optimum is open.

\paragraph{Adaptive refinement policies.}
Proposition~\ref{prop:refinement-helps} compares ``solve $h$ directly'' against ``refine $h$ then solve
its children'' as if they were exclusive. In practice, the executor first
attempts $h$ for some number of tries and refines only on persistent failure.
The right object is a sequential policy that, after $k$ failed attempts on
$h$, decides among continue, refine, or abandon, given a posterior over
$p_h$ and the remaining budget. Casting this as an optimal stopping problem
and characterising the resulting policy (presumably a threshold rule on the
posterior mean of $p_h$ relative to the per-child rates) is open.

\paragraph{Strengthening the contamination bound under correlated validator
errors.} The per-invocation union bound
$\bar N\,\alpha_{\max}^{\rightarrow}$ of
Theorem~\ref{thm:contamination-noisy-invocation} is loose when validator false
accepts are correlated across invocations, e.g.\ a learned validator with a
systematic blind spot for a particular invalid pattern. In that regime the
right object is a mixture model over latent ``vulnerability classes'' that
the validator either does or does not catch, with concentration measured
against the structure of the proposal distribution rather than the number of
attempts. We do not address this here.

\paragraph{Joint optimisation of validator and proposer.}
The framework treats $\LLM$ and $g$ as independent. In practice they share a
model family, and improving the validator can degrade the proposer (and vice
versa) under a fixed compute budget. The Pareto frontier between
$p_h$, $\alpha_h^{\rightarrow}$, and $\beta_h$ at fixed compute is the
practically relevant object, and is unaddressed.

\paragraph{Quantifying the latent-mode variance.}
Proposition~\ref{prop:correlation-upper-bound} reduces the gap between
realised and ideal retry amplification to the variance of $\rho(Z)$. Empirical
calibration of this variance for representative model/prompt families,
together with a quantitative theory of how diversity-inducing interventions
shrink it, would convert Theorem~\ref{thm:retry-amplification} from an asymptotic target into a usable
operating bound.

\section{Summary of Claims}

The transactional structured reasoning framework supports two core
theoretical claims, one structural and one stochastic.

\paragraph{Contamination safety and conditional correctness.}
Under clean initialization, exclusive commit authority, faithful rollback,
sound validation, final-answer dependence only on committed state, and
answer-extraction soundness, two things follow:
\begin{enumerate}[label=(\roman*)]
\item every committed state reachable by the executor is clean (Theorem~\ref{thm:contamination-safety});
\item if the executor emits a final answer, that answer is correct
(Corollary~\ref{cor:conditional-correctness}).
\end{enumerate}
Contamination safety is the structural part — bad samples don't enter
committed state — and is unconditional given the assumptions. Correctness
of an emitted answer additionally requires that the validator certify the
root certificate before extraction. Neither result claims the executor will
emit; failure modes (timeout, stall, budget exhaustion) remain.

\paragraph{Safe retry plus conditional recovery improvement.}
Treating the LLM as a stochastic proposal kernel, transactional retries
under exact validation cannot increase contamination risk relative to a
one-shot baseline (rejected attempts roll back to a clean state). The
quantitative improvement they provide depends on a local-rate model:
\begin{itemize}[leftmargin=*]
\item under conditional independence, the per-obligation rate amplifies from
$p_h$ to $1-(1-p_h)^{b_h}$ (Theorem~\ref{thm:retry-amplification});
\item under shared latent failure modes, that i.i.d.\ rate is an \emph{upper}
bound on what is achievable, and the realised rate can be as low as $p_h$
(Proposition~\ref{prop:correlation-upper-bound});
\item under a calibrated uniform persistence parameter $\rho_h$, the rate
is at least $1-(1-p_h)\rho_h^{\,b_h-1}$, interpolating between the two
extremes.
\end{itemize}
The global product lower bound (Theorem~\ref{thm:product-abstract}) is
stated abstractly in terms of per-obligation lower bounds
$q_h^{\mathrm{LB}}$, which the local-rate model supplies. Refinement is
beneficial when the analogous expression for refinement exceeds $q_h$ for
the same calibrated rates.

\paragraph{Scope.}
These are conditional guarantees, not claims that arbitrary long-horizon
reasoning is solved. The structural guarantees (no contamination, conditional
correctness) hold whenever the assumptions hold. The recovery improvement is
guaranteed to be safe but its magnitude is a function of validator strength,
local proposal quality, and sample diversity — none of which the framework
itself produces. The substantive engineering questions (constructing sound
validators for the task class, inducing enough sample diversity to push
$\rho_h$ below 1, allocating retry budget across obligations) remain open
and are surveyed in Section~\ref{sec:open-problems}.

\end{document}
